%%=============================================================================
%% Voorwoord
%%=============================================================================

\chapter*{\IfLanguageName{dutch}{Woord vooraf}{Preface}}%
\label{ch:voorwoord}

%% TODO:
%% Het voorwoord is het enige deel van de bachelorproef waar je vanuit je
%% eigen standpunt (``ik-vorm'') mag schrijven. Je kan hier bv. motiveren
%% waarom jij het onderwerp wil bespreken.
%% Vergeet ook niet te bedanken wie je geholpen/gesteund/... heeft

In dit voorwoord wil ik mijn motivatie voor het gekozen onderwerp toelichten en mijn dank uitspreken aan de mensen die mij tijdens deze bachelorproef hebben ondersteund.
Toen ik dit onderwerp in de lijst zag staan, sprak het mij meteen aan. De combinatie van artificiële intelligentie en een praktische toepassing in de retailsector sluit goed aan bij mijn interesses. 
Vooral het idee om computer vision toe te passen in een realistische omgeving, zoals een supermarkt, maakte dit onderwerp voor mij extra boeiend.
Tijdens het uitwerken van deze bachelorproef heb ik niet alleen mijn technische kennis verder kunnen verdiepen, maar ook inzicht gekregen in de uitdagingen die komen kijken bij het toepassen van AI in de praktijk. Dit maakte het proces zowel leerrijk als uitdagend.
Daarnaast wil ik graag mijn promotor bedanken voor de begeleiding en feedback tijdens het traject. Ook wil ik mijn co-promotor bedanken voor zijn steun gedurende het hele proces.
